\chapter{Transformation into SMT}
\label{ch:trans}

This chapter walks over the formal specification of the transformation process from \rego into \smtlib.
The structure of this chapter will correspond to the chapter \ref{ch:form}, which introduced the formal definition of targeted subset of the Rego policy language.

%%%%%%%%%%%%%%%%%%%%%%%%%%%%%%%%%%%%%%%%%%%%%%%%%%%%%%%%%%%%%%%
% TODO:
% variables
% values
%% atomic
%% composite
% references
% types 
% function calls
\section{Terms}
%%=============================================================

This section describes the transformation of Rego terms into SMT.
By the specification given in \ref{ch:form}, terms can be either values, variables, references, or function calls.

%%=============================================================
\subsection{Variables}
%%-------------------------------------------------------------

In SMT, variables can be declared in the global scope with the \texttt{declare-fun} construct.
Since \smtlib has is a strongly-typed language, we need to provide sorts of variables during the declaration.
For $x \in \Vars$ we set a its sort to be \texttt{OTypeD<N>}, where $N$ is the depth of its type, i.e., $\depth(x)$.
We will cover the specifics of such type later.

Intuitively, the global scope will only contain global variables.
Hypothetically, for no clash in variable names (which we ensure), also local variables could be declared with mentioned \texttt{declare-fun} construct.
We will choose to only specify them locally as expected by the Rego semantics but with the main purpose not to let them pollute the output model.
For this cause, we will use the \texttt{let} construct for assigned variables, \texttt{exists} for existentially quantified variables, and \texttt{forall} for universally quantified variables.








